\documentclass{beamer}
\usetheme{Madrid}
\usepackage{amsmath, amssymb, amsthm}
\usepackage{graphicx}
\usepackage{listings}
\usepackage{gensymb}
\usepackage[utf8]{inputenc}
\usepackage{hyperref}
\usepackage{gvv}
\usepackage{tikz}
\lstset{
  language=Python,
  basicstyle=\ttfamily\small,
  keywordstyle=\color{blue},
  stringstyle=\color{orange},
  numbers=left,
  numberstyle=\tiny\color{gray},
  breaklines=true,
  showstringspaces=false
}
\usetikzlibrary{decorations.pathmorphing}

\title{Question-5.8.20}
\author{EE25BTECH11022 - sankeerthan}
\date{}
\begin{document}

\frame{\titlepage}

\begin{frame}
\frametitle{Question}
Five years ago, Minu was thrice as old as Sonu.Ten years later,Minu will be twice as old as Sonu.How old are Minu and Sonu?
\end{frame}

\begin{frame}
\frametitle{Solution}
Let Minu's present age be x and Sonu's present age be y.
Five years ago, Minu was thrice as old as Sonu:
\begin{align}
\myvec{1 & 0}\myvec{x \\ y}-5 = 3\brak{\myvec{0 & 1}\myvec{x \\ y}-5} \\
\brak{\myvec{1 & 0}-3\myvec{0 & 1}}\myvec{x \\ y}= -10 \\ 
\myvec{1 & -3}\myvec{x \\ y}= -10 
\end{align}
Ten years later,Minu will be twice as old as Sonu:
\begin{align}
\myvec{1 & 0}\myvec{x \\ y}+10= 2\brak{\myvec{0 & 1}\myvec{x \\ y}+10} \\
\brak{\myvec{1 & 0}-2\myvec{0 & 1}}\myvec{x \\ y}= 10 \\
\myvec{1 & -2}\myvec{x \\ y}= 10 
\end{align}
\end{frame}

\begin{frame}
\frametitle{Solution}
Matrix form:
\begin{align}
\myvec{1 & -3 \\1 & -2}\myvec{x \\ y}=\myvec{-10 \\ 10}
\end{align}
Let 
\begin{align}
\vec{A} = \myvec{1 & -3\\ 1 & -2},\vec{X} =\myvec{x \\ y},\vec{B} =\myvec{-10 \\ 10}
\end{align}
Inverse of $\vec{A}$:
\begin{align}
\vec{A^{-1}} =\myvec{-2 & 3\\ -1 & 1} 
\end{align}
\end{frame}
\begin{frame}
\frametitle{Solution}
\begin{align}
\vec{X} &= \vec{A^{-1}}\vec{B} \\
&=\myvec{-2 & 3\\ -1 & 1}\myvec{-10 \\ 10} \\
&=\myvec{50 \\ 20}
\end{align}
Therefore, Minu's present age = 50 and Sonu's present age = 20
\end{frame}
\begin{frame}[fragile]
\frametitle{C-Code}
\begin{lstlisting}[language=C]
// code.c
#include <stdio.h>

void solve_ages(double *minu, double *sonu) {
    double A[2][2] = {{1, -3}, {1, -2}};
    double B[2] = {-10, 10};
    double det = A[0][0]*A[1][1] - A[1][0]*A[0][1];

    // Inverse matrix elements
    double invA[2][2] = {{A[1][1]/det, -A[0][1]/det},
                         {-A[1][0]/det, A[0][0]/det}};

    // Multiply invA by B
    *minu = invA[0][0]*B[0] + invA[0][1]*B[1];
    *sonu = invA[1][0]*B[0] + invA[1][1]*B[1];
}
\end{lstlisting}
\end{frame}
\begin{frame}[fragile]
\frametitle{Python-Code}
\begin{lstlisting}
# Code by GVV Sharma
# Modified for Problem Solution
# Released under GNU GPL
# Calculating area enclosed between curves
# call.py
import ctypes
import plot

lib = ctypes.CDLL('./code.so')
\end{lstlisting}
\end{frame}
\begin{frame}[fragile]
\frametitle{Python-Code}
\begin{lstlisting}
minu = ctypes.c_double()
sonu = ctypes.c_double()

lib.solve_ages(ctypes.byref(minu), ctypes.byref(sonu))

print(f"Minu's age: {minu.value}")
print(f"Sonu's age: {sonu.value}")

# Pass Python values to plot script
plot.plot_ages(minu.value, sonu.value)
\end{lstlisting}
\end{frame}

\end{document}